\section{Efficiency and parameter constraints}

This section gives symbolic efficiency formulas.  It does not certify concrete security.  A real submission should include a parameter estimator and independent cryptanalysis.

\subsection{Ciphertext and key sizes}

Let $N=n+1$ and $w_0=(k+1)(k+2)$.

\begin{center}
\begin{tabular}{lll}
\toprule
Object & Dense size & Sparse/structured size\\
\midrule
Compact input ciphertext & $N$ field elements & $O(k\log n+\log q)$ bits plus values\\
One GSW matrix ciphertext & $N^2$ field elements & $O(Nw)$ field/index pairs\\
Expansion key & $N$ GSW matrices & $O(N^2(k+2))$ field/index pairs\\
Compaction key & $N$ CLPN ciphertexts & $O(Nm_cn_c)$ dense, seedable in practice\\
Compacted output & $R$ CLPN ciphertexts & independent of circuit size\\
\bottomrule
\end{tabular}
\end{center}

The expansion key is large but amortized.  It is published once per key and reused for many input ciphertexts.  The final compact output does not grow with the evaluated circuit.

\subsection{Evaluation costs}

For row-sparse matrices with row supports $w_1,w_2$:
\[
    \EvalAdd: O(N(w_1+w_2)),
    \qquad
    \EvalMul: O(Nw_1w_2)
\]
field operations, ignoring sorting/merging overhead.  After depth $D$,
\[
    w_D=w_0^{2^D}.
\]
Thus \sable{} is a low-depth scheme.  It is realistic to target degree-2, degree-4, or possibly degree-8 computations before considering further compression or bootstrapping.

\subsection{Correctness budget}

A conservative correctness target is
\[
    \varepsilon_f+\varepsilon_{\rm cmp}(B)<\frac12.
\]
A stronger practical target is
\[
    \varepsilon_f+\varepsilon_{\rm cmp}(B)\le 0.1,
\]
so that a moderate replica count $R$ gives negligible final failure.  Using the bound from Section~6,
\[
    \varepsilon_f
    \le
    A(k+2)\eta\prod_{r=0}^{D-1}(1+w_0^{2^r}).
\]
Therefore a simple design inequality is
\[
    \eta
    \le
    \frac{\tau}{A(k+2)\prod_{r=0}^{D-1}(1+w_0^{2^r})}
\]
for target sparse-LPN evaluation error $\tau$.

For the compactor, if the final row support is $B$, the effective q-ary error rate is
\[
    \eta_{c,B}=\frac{q-1}{q}\left(1-\left(1-\frac{q\eta_c}{q-1}\right)^B\right).
\]
The code should be chosen so that decoding at this error rate fails with probability at most the desired $\varepsilon_{\rm cmp}(B)$.

\subsection{Illustrative non-secure toy parameters}

The following parameters are only for debugging a Python prototype:
\[
    q=127,\quad n=64,\quad k=3,\quad D=1,\quad R=9.
\]
They are intentionally small and not secure.  Their purpose is to validate algebraic correctness and serialization.

A research-grade estimator must choose $n,k,q,\eta,n_c,m_c,\eta_c$ jointly.  The estimator should output:
\begin{itemize}
    \item estimated classical and quantum cost of sparse-LPN distinguishing;
    \item estimated cost of dense q-ary LPN/code distinguishing;
    \item estimated ISD cost for the relevant code parameters;
    \item total number of LPN samples exposed by input ciphertexts and evaluation keys;
    \item correctness failure estimates before and after replication.
\end{itemize}

\subsection{Parameter-selection recipe}

A practical research workflow is:
\begin{enumerate}
    \item Fix the application class, especially multiplicative depth $D$ and addition budget $A$.
    \item Choose a small sparse row weight $k$ that keeps multiplication feasible.
    \item Compute $w_0=(k+1)(k+2)$ and the depth growth $w_D=w_0^{2^D}$.
    \item Choose $\eta$ small enough for correctness but not so small that low-noise LPN attacks become cheap.
    \item Choose $q$ to balance message space, q-ary error behavior, and implementation simplicity.
    \item Choose $\mathsf{Code}$ and $(n_c,m_c,\eta_c)$ so that compaction decoding succeeds at row support $B\le\min(N,w_D)$.
    \item Run attack estimators and increase dimensions until the desired security category is reached.
\end{enumerate}

\subsection{Benchmark baselines}

The natural baselines are TFHE/FHEW for Boolean circuits, BFV/BGV for exact arithmetic, and CKKS for approximate arithmetic \cite{DucasMicciancio2015,CGGI2020,FanVercauteren2012,BGV2014,CKKS2017}.  \sable{} should not be expected to dominate those mature lattice schemes broadly.  It should be evaluated on the specific workloads where code-based assumption diversity and low-degree structure matter.
